\documentclass[11pt,a4paper]{article}

\usepackage[T1]{fontenc}
\usepackage[english]{babel}
\usepackage{lmodern}
\usepackage{microtype}
\usepackage[a4paper,margin=25mm]{geometry}
\usepackage{amsmath,amssymb,amsthm,mathtools}
\usepackage{booktabs,tabularx,array}
\usepackage{xcolor}
\usepackage{enumitem}
\usepackage{fancyhdr}
\usepackage{hyperref}

\definecolor{navy}{HTML}{153A5B}
\definecolor{pale}{HTML}{EDF4F8}
\hypersetup{
  colorlinks=true,
  linkcolor=navy,
  citecolor=navy,
  urlcolor=navy,
  pdftitle={Exact Certificates for the (72,108) Frontier in the Two-Dimensional Jacobian Problem},
  pdfauthor={Billel Helali},
  pdfsubject={Exact algebraic certificates for the planar Jacobian problem},
  pdfkeywords={Jacobian conjecture, Keller map, computer-assisted proof, exact certificate, Newton polygon}
}

\pagestyle{fancy}
\fancyhf{}
\lhead{Exact exclusion of the $(72,108)$ frontier}
\rhead{B. Helali}
\cfoot{\thepage}
\setlength{\headheight}{14pt}

\newtheorem{theorem}{Theorem}[section]
\newtheorem{proposition}[theorem]{Proposition}
\newtheorem{lemma}[theorem]{Lemma}
\newtheorem{corollary}[theorem]{Corollary}
\theoremstyle{remark}
\newtheorem{remark}[theorem]{Remark}

\newcommand{\Q}{\mathbb{Q}}
\newcommand{\C}{\mathbb{C}}
\newcommand{\Jac}{\operatorname{Jac}}
\newcommand{\JC}{\mathrm{JC}}

\title{\textbf{Exact Certificates for the $(72,108)$ Frontier\\in the Two-Dimensional Jacobian Problem}}
\author{Billel Helali\\\small Independent Researcher, France}
\date{Version 1.1.0 preprint draft --- 21 July 2026}

\begin{document}
\maketitle

\begin{center}
\small Archival artifacts: \href{https://doi.org/10.5281/zenodo.21479814}{doi:10.5281/zenodo.21479814}
\end{center}

\begin{center}
\fcolorbox{navy}{pale}{%
\begin{minipage}{0.91\textwidth}
\textbf{Status and scope.}
The identities reported below have been replayed in exact characteristic-zero arithmetic for the transcribed Laurent/Newton systems. The exclusion of the degree pair $(72,108)$, and hence the degree-$125$ lower bound, remains conditional on the correctness and exhaustiveness of the published reduction of Guccione--Guccione--Horruitiner--Valqui and on its faithful transcription. This note does \emph{not} prove the two-dimensional Jacobian conjecture. External mathematical and computational review is required before the result should be treated as established.
\end{minipage}}
\end{center}

\begin{abstract}
Guccione, Guccione, Horruitiner and Valqui proved that a hypothetical two-dimensional Jacobian counterexample of maximum degree below $125$ must have degree pair $(72,108)$, up to order, and reduced that pair to two Laurent/Newton configurations. We give exact characteristic-zero ideal certificates for the complete coefficient systems transcribed from those configurations. The hard certificate is an $89.1$ MB identity over a degree-five number field, reconstructed from a fixed $1925\times1925$ rational subsystem and verified against all $2010$ scalar equations. A separate \texttt{gmpy2} implementation checks $13{,}410$ coefficient-field products, equivalent to $335{,}250$ scalar products. We also provide a source-to-system audit, a clean-room Linux replay from the frozen Zenodo deposit, and deliberate corruption tests. The exact identities exclude the transcribed systems. The consequence that $(72,108)$ itself is impossible, and hence that any planar counterexample has maximum degree at least $125$, remains conditional on the correctness and exhaustiveness of the published reduction and on external confirmation of our transcription.
\end{abstract}

\tableofcontents

\section{Introduction and scope}

For a polynomial map $F=(F_1,\ldots,F_n):\C^n\to\C^n$, the classical Jacobian conjecture asserts that a nonzero constant determinant $\det JF$ forces $F$ to have a polynomial inverse. The conjecture was introduced by Keller in 1939 \cite{Keller1939}. A three-dimensional counterexample was publicly announced in July 2026 \cite{Alpoge2026}; that development provides current context but plays no role in the argument below. This paper concerns only the still-separate plane problem.

Guccione, Guccione, Horruitiner and Valqui (GGHV) showed that every hypothetical plane counterexample of maximum degree below $125$ must have degree pair $(72,108)$, up to order \cite{GGHV2022}. Their Proposition~4.3 reduces the associated $(8,28)$ corner case to two explicit Laurent/Newton configurations. Our contribution is the reconstruction and exact replay of algebraic certificates excluding the coefficient systems transcribed from both configurations.

The central logical distinction is essential. Our verifiers prove polynomial identities for explicit systems. The inference from those systems to all maps of degree pair $(72,108)$ additionally uses GGHV's published reduction and our interface to it. We audit that interface below and release an executable checker, but an external specialist review is still pending. Accordingly, the unconditional theorem in this paper is the exclusion of the two transcribed systems; the degree-$125$ statement is a conditional corollary.

\section{The planar \texorpdfstring{$(72,108)$}{(72,108)} frontier}

The main published input is the following consequence of Theorem~2.1 and Proposition~4.3 of \cite{GGHV2022}: below maximum degree $125$, only $(72,108)$ and its symmetric pair remain, and a counterexample of that degree type would yield Laurent polynomials $P,Q\in K[x,x^{-1},y]$ over an algebraically closed characteristic-zero field satisfying
\[
[P,Q]=x^2
\]
in one of the following two Newton configurations:
\begin{align*}
\text{Case 1:}\quad
N(P)&=\operatorname{conv}\{(0,0),(1,0),(8,14),(8,16),(0,8)\},\\
N(Q)&=\operatorname{conv}\{(0,0),(2,1),(12,21),(12,24),(0,12)\};\\[0.4em]
\text{Case 2:}\quad
N(P)&=\operatorname{conv}\{(0,0),(1,0),(8,14),(8,16)\},\\
N(Q)&=\operatorname{conv}\{(0,0),(2,1),(12,21),(12,24)\}.
\end{align*}

We emphasize the logical boundary: our computation proves identities for the systems obtained from these two polygonal cases. Promoting those identities to a theorem about all plane maps requires the reduction in \cite{GGHV2022} to be exhaustive and the local transcription and normalization to be exact.

\section{Audit of the published-to-implemented interface}

The primary-source audit uses arXiv:2204.14178v1. Theorem~2.1 appears on PDF page~2, and Proposition~4.3 with its proof appears on PDF pages~10--12. The ring $L^{(1)}$, the bracket $[P,Q]=x^2$, and all ten displayed vertices above agree with those pages. SHA-256 digests of the downloaded PDF and TeX source, together with a line-by-line claim ledger, are recorded in \texttt{TRANSCRIPTION\_AUDIT.md}.

\subsection{Lattice supports}

Under $t=xy^2$ and $z=y^{-1}$, a Laurent monomial transforms as
\[
x^iy^j=t^iz^{2i-j}.
\]
Fresh boundary-inclusive lattice enumeration gives $61$ and $125$ points for the two Case~1 polygons, and $25$ and $47$ points for Case~2. The corresponding $z$-bands are
\[
\begin{array}{c|cc}
&P&Q\\ \hline
\text{Case 1}&-8,-7,\ldots,2&-12,-11,\ldots,3\\
\text{Case 2}&0,1,2&0,1,2,3.
\end{array}
\]
Thus the Case~2 ansatz in the next section includes every lattice-supported coefficient. Case~1 begins with the same upper bands and introduces complete descending slices. A hypothetical full Case~1 solution must satisfy the first five descending bracket layers used in the computation, so inconsistency of that necessary subsystem is sufficient.

Constants are removed from $P$ and $Q$ by translation in the target; the bracket is unchanged. This operation maps every pair with the published polygons into the implemented ansatz, which is the implication required for an exclusion.

\subsection{Reversible vertex normalization}

Let $a_1,a_8,c_8\ne0$ be the coefficients of $P$ at $(1,0),(8,14),(8,16)$. For
\[
\widetilde P=\rho P(\lambda x,\mu y),\qquad
\widetilde Q=\sigma Q(\lambda x,\mu y),
\]
choose
\[
\mu^2=\frac{a_8}{c_8},\qquad
\lambda^7=\frac{a_1}{a_8}\mu^{-14},\qquad
\rho=(a_1\lambda)^{-1},\qquad
\sigma=(\rho\lambda^3\mu)^{-1}.
\]
The algebraically closed ground field supplies the required nonzero roots. The three selected coefficients of $\widetilde P$ equal one. The chain rule gives
\[
[\widetilde P,\widetilde Q]=\rho\sigma\lambda^3\mu x^2=x^2.
\]
This proves that the assignments $a_1=a_8=c_8=1$ do not discard a component. The released bridge checker verifies these scaling identities symbolically.

\subsection{Division and saturation ledger}

All pivots used by the exact linear solves are nonzero elements of a number field, not parameter polynomials. Compatibility rows are retained. On each Case~1 sign branch, the elimination of $r$ follows from an identity $E_6-\lambda E_2=S\Lambda^2$, with $S\ne0$ and $\Lambda$ monic affine-linear in $r$; hence $\Lambda=0$ at every common zero. Equations removed afterward are checked to be nonzero field-scalar multiples of retained equations.

The only variable localization occurs when $u_3=N/h$ is used in the declared branch $h\ne0$. The complementary specialization $h=0$ is verified against the original seven pre-division equations by separate unit certificates whose payloads contain hashes of those source systems. The regenerated $h\ne0$ equations contain no additional common power of $h$ to divide out. Consequently, the local computation contains no unrecorded saturation.

The executable audit \texttt{scripts/verify\_transcription\_bridge.py} independently checks the polygon bands, normalization, coefficient equations, exhaustive sign split, affine eliminations, proportionality removals, regenerated branch systems, and the $h=0/h\ne0$ coverage. Its scope starts at the statement of Proposition~4.3; it does not re-prove that proposition or its cited predecessors.

\section{Exact Laurent reduction}

Set $t=xy^2$ and $z=y^{-1}$. Then $[t,z]_{x,y}=-1$ and $x^2=t^2z^4$. The relevant upper bands can be written
\[
P=A(t)z^2+B(t)z+C(t),\qquad
Q=D(t)z^3+E(t)z^2+F(t)z+G(t).
\]
Equating coefficients in $[P,Q]=x^2$ gives
\begin{align*}
2AD'-3A'D&=t^2, \tag{J4}\label{eq:J4}\\
2(AE'-A'E)+(BD'-3B'D)&=0, \tag{J3}\label{eq:J3}\\
(2AF'-A'F)+(BE'-2B'E)-3C'D&=0, \tag{J2}\label{eq:J2}\\
2AG'+(BF'-B'F)-2C'E&=0, \tag{J1}\label{eq:J1}\\
BG'-C'F&=0. \tag{J0}\label{eq:J0}
\end{align*}
The replay checks the coordinate bracket, the identity $x^2=t^2z^4$, all five formulas, and the polygonal lattice supports exactly.

With normalized nonzero vertex coefficients,
\[
A=t+a_2t^2+\cdots+a_7t^7+t^8,
\qquad D=d_2t^2+\cdots+d_{12}t^{12}.
\]
Equation~\eqref{eq:J4} solves the eleven coefficients of $D$ and leaves six compatibility equations in $a_2,\ldots,a_7$. Exact rational Gr\"obner and FGLM calculations yield an irreducible degree-$35$ polynomial $H(a_7)$ together with five relations linear in $a_2,\ldots,a_6$. Subsequent residual systems descend to
\[
L=\Q[w]/(w^5-w^4+3w^3+3w^2+26).
\]

\section{Certificate architecture}

Two elementary observations carry the logical proof. If $1=\sum_i T_iE_i$, then the $E_i$ have no common zero. If $h=\sum_iT_iE_i$, then every common zero satisfies $h=0$. All displayed certificates are identities in polynomial rings over exact number fields.

\subsection{Case 2}

Exact elimination in equations~\eqref{eq:J3}--\eqref{eq:J0} leaves four residual generators $R_1,\ldots,R_4$ over a degree-$35$ field. A Macaulay system of degree $8$, with dimensions $165\times182$, produces exact multipliers satisfying
\[
1=T_1R_1+T_2R_2+T_3R_3+T_4R_4.
\]
The serialized certificate is evaluated again over the number field; the exact remainder is $1$. Hence Case~2 has no solution.

\subsection{Case 1: exhaustive branching and the specialization \texorpdfstring{$h=0$}{h=0}}

The negative $z$-bands produce thirteen compatibility equations. Their first factorization forces the exhaustive split $s=c$ or $s=-c$, where $c\neq0$. In each branch a checked relation
\[
E_6-\lambda E_2=S\Lambda^2,\qquad S\neq0,
\]
forces an affine-linear form $\Lambda$ to vanish, eliminating one variable without division by a variable expression.

For $h=0$, exact unit-ideal certificates are reconstructed separately for the two branches using the seven pre-division equations:
\[
1=\sum_i A_iE_i\big|_{h=0}.
\]
Thus neither branch has a solution with $h=0$. The use of pre-division equations is important: it avoids losing a component through an unrecorded saturation or variable division.

\subsection{Case 1: the hard identity}

For the branch $s=c$, four reduced generators in $L[h,u_1,u_2]$ admit the identity
\begin{equation}\label{eq:hard}
\boxed{h=T_1E_1+T_2E_2+T_3E_3+T_4E_4.}
\end{equation}
Therefore every common zero would satisfy $h=0$, which the preceding unit certificate excludes. The involution
\[
(h,u_1,u_2)\longmapsto(h,-u_1,-u_2)
\]
transports the equations and identity~\eqref{eq:hard} exactly to $s=-c$. Both branches, and therefore Case~1, are empty.

\section{Construction of the hard certificate}

The multiplier ansatz for~\eqref{eq:hard} has $385$ coefficients in $L$ and gives $402$ coefficient equations in $L$. Expanding on the basis $1,w,\ldots,w^4$ produces a rational linear system with
\[
2010\ \text{equations and}\ 1925\ \text{unknowns}.
\]
A deterministic set of $1925$ scalar rows was selected by exact elimination modulo $71$. The square minor $B_{71}$ satisfies
\[
\det(B_{71})\equiv10\pmod{71},
\]
so the chosen support is nonsingular modulo $71$ and hence nonsingular over $\Q$. The corresponding $1925\times1925$ system was rebuilt and solved in exact rational arithmetic using FLINT. The resulting vector has $1925$ nonzero rational coordinates; its largest numerator and denominator have $85{,}342$ and $85{,}275$ bits, respectively. After serialization, the human-readable identity occupies $89{,}105{,}967$ bytes.

The selected square subsystem is only the construction mechanism. Soundness comes from substituting the recovered multipliers into all $402$ equations over $L$, equivalently all $2010$ rational scalar rows, and checking a zero residual everywhere.

\section{Certificate formats and trust boundaries}

The three smaller certificates are Python pickle payloads with explicit format identifiers, minimal-polynomial coefficients, serialized generators and serialized multipliers. The Case~2 verifier regenerates its selected generators from the coefficient equations and requires exact equality with the embedded copy before evaluating the unit identity. Each $h=0$ payload records the file name and SHA-256 of its pre-division source system; the verifier regenerates the specialization and compares every term.

The hard certificate is a line-oriented text file. Its four-line header records the degree-five minimal polynomial, the identity and the number of multiplier coefficients. Each subsequent record has the form
\begin{verbatim}
C|equation|h_exp|u1_exp|u2_exp|field_coefficient
\end{verbatim}
where a field coefficient is a reduced element of $\Q[w]/(w^5-w^4+3w^3+3w^2+26)$. The independent parser rejects malformed lines, duplicate supports, unreduced powers and an incorrect header. It then constructs the full rational coefficient map from the separately serialized generators.

Pickle is used only as a data encoding and should be opened solely from the hashed archive in a disposable environment. It is not a security boundary. Mathematical soundness comes from generator regeneration and exact identity evaluation, not from trusting serialized claims.

\section{Independent replay and reproducibility}

The full frozen replay is launched from the exact workspace by
\begin{verbatim}
./verify_all.sh
\end{verbatim}
and terminates with
\begin{verbatim}
ALL_SERIALIZED_EXACT_CERTIFICATES_PASS
GMPY2_EXACT_PASS
SYSTEM_SYMMETRY_PASS
BRANCH1_EXACT_IDENTITY_PASS
BRANCH2_EXACT_IDENTITY_PASS
JC2_72_108_EXACT_REPLAY_PASS
\end{verbatim}

The independent \texttt{gmpy2} path parses the large text certificate, normalizes its $1925$ rational coordinates, and evaluates the identity without relying on the FLINT solver that produced it. It performs $13{,}410$ coefficient-field products, corresponding to $335{,}250$ scalar products. The branch symmetry is also evaluated as a polynomial identity rather than inferred from matching dimensions.

For clean-room reproduction, the repository supplies a Docker image based on CPython 3.14.6. The base image is pinned by manifest digest. The packages \texttt{gmpy2} 2.3.1, \texttt{numpy} 2.5.1, \texttt{python-flint} 0.9.0, \texttt{sympy} 1.14.0 and \texttt{mpmath} 1.3.0 are pinned with Linux-wheel hashes. The container downloads the frozen $v1.0.1$ snapshot from Zenodo only, verifies the outer and nested SHA-256 digests, and then runs the exact replay, bridge audit and negative controls. The corresponding GitHub Actions job uses Ubuntu 24.04 on x86-64 and requires the final marker
\begin{verbatim}
JC2_72_108_CLEAN_ROOM_PASS
\end{verbatim}
before success.

Five negative controls deliberately change, in turn, a source hash, a validly serialized certificate coefficient, an equation coefficient, a branch-symmetry coefficient and one exact rational coordinate of the hard certificate. Every altered object is required to fail its intended verifier. The tests operate on temporary copies and leave the frozen archive untouched.

\begin{table}[ht]
\centering
\caption{Claim-to-verifier summary.}
\small
\begin{tabularx}{\textwidth}{@{}>{\raggedright\arraybackslash}p{0.25\textwidth}>{\raggedright\arraybackslash}p{0.29\textwidth}>{\raggedright\arraybackslash}X@{}}
\toprule
Claim & Exact object & Verifier \\
\midrule
Case 2 is empty & $1=\sum T_iR_i$ & generator regeneration and serialized-certificate replay \\
Case 1 split is exhaustive & $\kappa(s^2-c^2)^2$ & transcription-bridge checker \\
$r$ elimination is component-preserving & $E_6-\lambda E_2=S\Lambda^2$ & transcription-bridge checker \\
The $h=0$ branches are empty & two pre-division unit identities & source-hash, generator and identity replay \\
The hard branch forces $h=0$ & $h=\sum T_iE_i$ & independent \texttt{gmpy2} scalar replay \\
The second sign branch is covered & exact involution & branch-symmetry verifier \\
The published degree consequence applies & GGHV Theorem~2.1 and Proposition~4.3 & external mathematical audit still required \\
\bottomrule
\end{tabularx}
\end{table}

\begin{table}[ht]
\centering
\caption{Core reconstructed artifacts. SHA-256 digests identify the exact files used in this replay.}
\scriptsize
\begin{tabularx}{\textwidth}{@{}>{\raggedright\arraybackslash}p{0.30\textwidth}>{\raggedright\arraybackslash}X@{}}
\toprule
Artifact & SHA-256 \\
\midrule
\texttt{hard/h\_certificate\_exact.txt} & \texttt{0e48ffab32469ef8405a6945b16cf1521ddeb3c592ae4e5051968110a4dc656a} \\
\texttt{case2\_exact\_certificate.pkl} & \texttt{cfbc3c39d7a28013671144f43ef76f0498542eaf6d562dd624bba3311194e4aa} \\
\texttt{h0\_branch1\_exact\_certificate.pkl} & \texttt{664de005e99bc6a0e61ba479ba64ed57ddbfa9c5399b955cbb12237ac70f8186} \\
\texttt{h0\_exact\_certificate.pkl} & \texttt{d7844c2f8edea62be4e3a7fe8a160dc4b7b70efd1262993ec6d2ed7e78722a1a} \\
\texttt{verify\_all.sh} & \texttt{f40bacfa84d915b375b4158109d5120b2d964a8d8012be770750d310abfb5837} \\
\bottomrule
\end{tabularx}
\end{table}

The reconstructed objects are mathematically equivalent to the historically described certificates but not byte-for-byte copies; serialization and textual formatting changed, so their hashes necessarily differ. The manifest \texttt{RECONSTRUCTED\_CERTIFICATES.sha256} is the authority for this replay.

\section{Main result and limitations}

\begin{theorem}[Exact exclusion of the transcribed systems]
The two characteristic-zero coefficient systems transcribed from Cases 1 and 2 of the $(72,108)$ Laurent/Newton reduction have no solutions.
\end{theorem}
\begin{proof}
Case~2 is inconsistent because its four residual generators admit an exact unit certificate. In Case~1 the split $s=\pm c$ is exhaustive. On each branch, the hard identity forces $h=0$, while the corresponding pre-division unit certificate excludes $h=0$. The exact involution verifies the transport from the first branch to the second.
\end{proof}

\begin{corollary}[Conditional degree bound]
Assume that Proposition~4.3 of \cite{GGHV2022} is exhaustive for the $(72,108)$ configuration and that the local normalization and transcription are faithful. Then $(72,108)$ cannot be the degree pair of a plane Jacobian counterexample. Consequently, every hypothetical counterexample to $\JC(2)$ has
\[
\max\{\deg P,\deg Q\}\geq125.
\]
\end{corollary}

This conclusion does not prove $\JC(2)$: configurations of maximum degree at least $125$ remain outside its scope. It also does not replace a line-by-line external audit of the passage from the published Newton polygons to the implemented coefficient systems. The appropriate public claim is therefore an \emph{exact computer-assisted candidate exclusion}, conditional on that theoretical interface, rather than a completed proof of the plane conjecture.

\section{Conclusion}

The exact reconstruction reported here excludes both coefficient systems transcribed from the two alternatives of GGHV Proposition~4.3. The computational content is finite and replayable: regenerated supports and residuals, three unit certificates, a large exact ideal-membership identity, fixed hashes, an independent arithmetic verifier, an executable interface audit and deliberate failure tests. Conditional on the published reduction and a faithful interface, this removes the sole degree pair left below $125$. The next necessary step is independent specialist review of the upstream reduction and of the transcription documented here; the conditional qualifier should remain until that review is obtained.

\section*{Data and code availability}
The LaTeX source, exact certificates, reconstruction programs, manifests, replay scripts and audit materials are archived at \href{https://doi.org/10.5281/zenodo.21479814}{doi:10.5281/zenodo.21479814} and developed at \url{https://github.com/bilLkarkariy/jc2-72-108-exact-certificates}. The principal certificate is too large to print but is directly machine-verifiable from \texttt{hard/h\_certificate\_exact.txt} in the frozen replay bundle. The files \texttt{TRANSCRIPTION\_AUDIT.md}, \texttt{CLAIMS\_TO\_FILES.md} and \texttt{REPRODUCE.md} state the trust boundary and the complete clean-room procedure.

\section*{Acknowledgment of computational assistance}
OpenAI Codex and GPT-5.6 Sol were used interactively to assist with symbolic derivations, software development, certificate reconstruction, verification design and manuscript preparation. Every central computational claim is supported by explicit characteristic-zero identities and deterministic replay programs. The human author reviewed the outputs and accepts full responsibility for the manuscript and released artifacts.

\begin{thebibliography}{9}
\bibitem{Keller1939}
O.-H. Keller,
\emph{Ganze Cremona-Transformationen},
Monatshefte f\"ur Mathematik und Physik \textbf{47} (1939), 299--306.

\bibitem{GGHV2022}
J. A. Guccione, J. J. Guccione, R. Horruitiner, and C. Valqui,
\emph{Increasing the degree of a possible counterexample to the Jacobian Conjecture from 100 to 108},
arXiv:2204.14178 (2022),
\url{https://arxiv.org/abs/2204.14178}.

\bibitem{Alpoge2026}
L. Alp\"oge,
\emph{Public announcement of an explicit three-dimensional Keller map},
X post, 20 July 2026,
\url{https://x.com/__alpoge__/status/2079028340955197566}.

\end{thebibliography}

\end{document}
